%============================================================
% FORMAT UNESA: hasil kalibrasi TEMPLATE FORMAT *.docx
% Jangan ubah tanpa panduan kampus baru atau persetujuan pembimbing.
%============================================================

\geometry{
  a5paper,
  twoside,
  asymmetric,
  left=2.5cm,
  right=2cm,
  top=2.5cm,
  bottom=2cm,
  headheight=12pt,
  headsep=0pt,
  footskip=0.9cm
}

\setstretch{1.15}
\setlength{\parindent}{1cm}
\setlength{\parskip}{0pt}
\setlength{\emergencystretch}{1.5em}
\raggedbottom

\definecolor{unesagray}{HTML}{5F5D5D}
\definecolor{unesagold}{HTML}{E8CE69}

% Judul bab resmi: BAB I lalu nama bab, Book Antiqua 10 pt tebal dan terpusat.
\renewcommand{\thechapter}{\Roman{chapter}}
\renewcommand{\thesection}{\Alph{section}}
\renewcommand{\thesubsection}{\thesection.\arabic{subsection}}
\renewcommand{\thesubsubsection}{\thesubsection.\arabic{subsubsection}}
\titleformat{\chapter}[display]
  {\centering\bfseries\fontsize{10}{13.8}\selectfont}
  {BAB \thechapter}
  {0pt}
  {}
\titlespacing*{\chapter}{0pt}{-18pt}{13.8pt}

% Subbagian template memakai A., B., C., dan seterusnya.
\titleformat{\section}
  {\bfseries\fontsize{10}{13.8}\selectfont}
  {\makebox[0.635cm][l]{\thesection.}}
  {0pt}
  {}
\titlespacing*{\section}{0.63cm}{13.8pt}{0pt}

\titleformat{\subsection}
  {\bfseries\fontsize{10}{13.8}\selectfont}
  {\thesubsection}
  {0.45em}
  {}
\titlespacing*{\subsection}{1.27cm}{6.9pt}{0pt}

\titleformat{\subsubsection}
  {\bfseries\fontsize{10}{13.8}\selectfont}
  {\thesubsubsection}
  {0.45em}
  {}
\titlespacing*{\subsubsection}{1.27cm}{6.9pt}{0pt}

\setcounter{secnumdepth}{3}
\setcounter{tocdepth}{3}

% Nomor tabel, gambar, dan persamaan tetap memakai nomor bab Arab.
\renewcommand{\thefigure}{\arabic{chapter}.\arabic{figure}}
\renewcommand{\thetable}{\arabic{chapter}.\arabic{table}}
\renewcommand{\theequation}{\arabic{chapter}.\arabic{equation}}

\captionsetup{
  font=normalsize,
  labelfont=normalfont,
  textfont=normalfont,
  justification=centering,
  singlelinecheck=false,
  labelsep=period,
  skip=4pt
}
\captionsetup[table]{position=above}
\renewcommand{\figurename}{Gambar}
\renewcommand{\tablename}{Tabel}

\setlength{\intextsep}{8pt}
\setlength{\textfloatsep}{8pt}
\setlength{\floatsep}{8pt}

\SetTblrInner{
  hlines,
  vlines,
  row{1}={font=\bfseries,halign=c,valign=m},
  rows={valign=m},
  rowsep=2pt,
  colsep=3pt
}

% Nomor halaman di bawah tengah, tanpa header dan garis.
\fancypagestyle{unesa}{
  \fancyhf{}
  \fancyfoot[C]{\thepage}
  \renewcommand{\headrulewidth}{0pt}
  \renewcommand{\footrulewidth}{0pt}
}
\fancypagestyle{plain}{
  \fancyhf{}
  \fancyfoot[C]{\thepage}
  \renewcommand{\headrulewidth}{0pt}
  \renewcommand{\footrulewidth}{0pt}
}
\pagestyle{unesa}

% Halaman kosong saat transisi ke halaman ganjil pada dokumen bolak-balik.
\makeatletter
\renewcommand{\cleardoublepage}{%
  \clearpage
  \if@twoside
    \ifodd\c@page\else
      \hbox{}\vfill
      \begin{center}(Halaman ini sengaja dikosongkan)\end{center}
      \vfill
      \thispagestyle{empty}
      \newpage
    \fi
  \fi
}
\makeatother

% Judul dan isi daftar otomatis.
\renewcommand{\contentsname}{DAFTAR ISI}
\renewcommand{\listtablename}{DAFTAR TABEL}
\renewcommand{\listfigurename}{DAFTAR GAMBAR}
\renewcommand{\bibname}{DAFTAR PUSTAKA}

\renewcommand{\cfttoctitlefont}{\hfill\bfseries}
\renewcommand{\cftaftertoctitle}{\hfill\null}
\renewcommand{\cftlottitlefont}{\hfill\bfseries}
\renewcommand{\cftafterlottitle}{\hfill\null}
\renewcommand{\cftloftitlefont}{\hfill\bfseries}
\renewcommand{\cftafterloftitle}{\hfill\null}
\setlength{\cftbeforetoctitleskip}{-18pt}
\setlength{\cftbeforelottitleskip}{-18pt}
\setlength{\cftbeforeloftitleskip}{-18pt}
\setlength{\cftaftertoctitleskip}{13.8pt}
\setlength{\cftafterlottitleskip}{13.8pt}
\setlength{\cftafterloftitleskip}{13.8pt}
\renewcommand{\cftdotsep}{1}
\setlength{\cftbeforechapskip}{0pt}
\setlength{\cftbeforesecskip}{0pt}
\setlength{\cftbeforesubsecskip}{0pt}
\renewcommand{\cftchappresnum}{BAB~}
\settowidth{\cftchapnumwidth}{BAB~VIII~~}
\renewcommand{\cftsecaftersnum}{.}
\renewcommand{\cftfigpresnum}{Gambar~}
\settowidth{\cftfignumwidth}{Gambar~10.10~~}
\renewcommand{\cfttabpresnum}{Tabel~}
\settowidth{\cfttabnumwidth}{Tabel~10.10~~}

\newlistof{lampiran}{loa}{DAFTAR LAMPIRAN}
\renewcommand{\cftloatitlefont}{\hfill\bfseries}
\renewcommand{\cftafterloatitle}{\hfill\null}
\setlength{\cftbeforeloatitleskip}{-18pt}
\setlength{\cftafterloatitleskip}{13.8pt}

\newcommand{\DaftarIsiUNESA}{%
  \clearpage\phantomsection
  \renewcommand{\contentsname}{DAFTAR ISI}%
  \addcontentsline{toc}{chapter}{DAFTAR ISI}
  \begingroup\singlespacing\sloppy\tableofcontents\endgroup
}
\newcommand{\DaftarTabelUNESA}{%
  \clearpage\phantomsection
  \renewcommand{\listtablename}{DAFTAR TABEL}%
  \addcontentsline{toc}{chapter}{DAFTAR TABEL}
  \begingroup\singlespacing\listoftables\endgroup
}
\newcommand{\DaftarGambarUNESA}{%
  \clearpage\phantomsection
  \renewcommand{\listfigurename}{DAFTAR GAMBAR}%
  \addcontentsline{toc}{chapter}{DAFTAR GAMBAR}
  \begingroup\singlespacing\listoffigures\endgroup
}
\newcommand{\DaftarLampiranUNESA}{%
  \clearpage\phantomsection
  \addcontentsline{toc}{chapter}{DAFTAR LAMPIRAN}
  \begingroup\singlespacing\listoflampiran\endgroup
}

\newcommand{\DaftarPustakaUNESA}{%
  \clearpage\phantomsection
  \renewcommand{\bibname}{DAFTAR PUSTAKA}%
  \addcontentsline{toc}{chapter}{DAFTAR PUSTAKA}
  \ifshowbibliography
    \bibliographystyle{apalike}
    \bibliography{refs/references}
  \else
    \chapter*{DAFTAR PUSTAKA}
    \noindent [Daftar pustaka belum diisi. Tambahkan hanya sumber yang telah
    diverifikasi dan benar-benar disitasi dalam naskah.]
  \fi
}

\newcommand{\FrontTitle}[2][]{%
  \clearpage
  \phantomsection
  \def\FrontTocTitle{#1}%
  \ifx\FrontTocTitle\empty
    \addcontentsline{toc}{chapter}{#2}%
  \else
    \addcontentsline{toc}{chapter}{#1}%
  \fi
  \begin{center}\bfseries #2\end{center}
  \vspace{13.8pt}
}

\newcommand{\IdentityRow}[2]{%
  \noindent\begin{tabularx}{\textwidth}{@{}p{3.05cm}@{\hspace{0.1cm}:\hspace{0.15cm}}>{\raggedright\arraybackslash}X@{}}
    #1 & #2
  \end{tabularx}\par
}

\newcommand{\AbstractIdentityRow}[2]{%
  \noindent\begin{tabularx}{\textwidth}{@{}p{3.8cm}@{\hspace{0.1cm}:\hspace{0.15cm}}>{\raggedright\arraybackslash}X@{}}
    #1 & #2
  \end{tabularx}\par
}

\newcommand{\CheckedBox}{$\boxtimes$}
\newcommand{\EmptyBox}{$\square$}

\newenvironment{kutipanpanjang}{%
  \begin{list}{}{%
    \setlength{\leftmargin}{1cm}%
    \setlength{\rightmargin}{1cm}%
    \setlength{\topsep}{6.9pt}%
    \setlength{\parsep}{0pt}%
  }\item\singlespacing\ignorespaces
}{\end{list}}

\newcommand{\Sumber}[1]{%
  \par\vspace{2pt}{\raggedright Sumber: #1\par}
}

\newcommand{\Lampiran}[1]{%
  \refstepcounter{lampiran}
  \chapter*{LAMPIRAN \thelampiran\\#1}
  \phantomsection
  \addcontentsline{toc}{chapter}{LAMPIRAN \thelampiran. #1}
  \addcontentsline{loa}{lampiran}{\protect\numberline{\thelampiran}#1}
}

\hypersetup{
  pdftitle={\JudulPenelitian},
  pdfauthor={\NamaMahasiswa},
  pdfsubject={\JenisDokumenCover\ Universitas Negeri Surabaya},
  pdfkeywords={\KataKunci}
}
\pdfstringdefDisableCommands{\def\\{ }}
